%% kapa: exam=1 round=32 year=2021 period=1 form=A subject=real_estate subject_ko=부동산학원론 seq=40 number=120
%% source: Q-Net 감정평가사 기출문제 (cst00309) artlSeq=5208397 file=extracted/1008/5208397/A형/2021년도 제32회 감정평가사 1차 1교시 A형.pdf
%% conversion: kordoc 4.24.51 -> tools/gen1.py (auto); review=proofread-sonnet; flags=none
\begin{question}{120}
감정평가 실무기준에서 규정하고 있는 수익환원법에 관한 내용으로 옳지 않은 것은?
\begin{choices}
  \item 수익환원법으로 감정평가할 때에는 직접환원법이나 할인현금흐름분석법 중에서 감정평가 목적이나 대상물건에 적절한 방법을 선택하여 적용한다.
  \item 부동산의 증권화와 관련한 감정평가 등 매기의 순수익을 예상해야 하는 경우에는 할인현금흐름분석법을 원칙으로 하고 직접환원법으로 합리성을 검토한다.
  \item 직접환원법에서 사용할 환원율은 요소구성법으로 구하는 것을 원칙으로 한다. 다만, 요소구성법의 적용이 적절하지 않은 때에는 시장추출법, 투자결합법, 유효총수익승수에 의한 결정방법, 시장에서 발표된 환원율 등을 검토하여 조정할 수 있다.
  \item 할인현금흐름분석법에서 사용할 할인율은 투자자조사법(지분할인율), 투자결합법(종합할인율), 시장에서 발표된 할인율 등을 고려하여 대상물건의 위험이 적절히 반영되도록 결정하되 추정된 현금흐름에 맞는 할인율을 적용한다.
  \item 복귀가액 산정을 위한 최종환원율은 환원율에 장기위험프리미엄·성장률·소비자물가 상승률 등을 고려하여 결정한다.
\end{choices}
\end{question}
\answer{3}
